% Source: https://artofproblemsolving.com/community/c3h1700210_math_is_this_possible
% Stable identity: aops-topic-1700210-q01
\textbf{Yes.} Take the ten numbers
\[
-2,\ 1,\ 1,\ 1,\ 1,\ 3,\ 3,\ 3,\ 3,\ 6 .
\]
\begin{itemize}
\item Mean: $\dfrac{-2+4\cdot1+4\cdot3+6}{10}=\dfrac{20}{10}=2$.
\item Median: the 5th and 6th values are $1$ and $3$, so the median is $2$.
\item Quartiles: the lower half $-2,1,1,1,1$ has median $Q_1=1$, the upper half $3,3,3,3,6$ has median $Q_3=3$, so the IQR is $3-1=2$.
\item Standard deviation: the deviations from the mean are $-4$, four times $-1$, four times $1$, and $4$, so the variance is $\dfrac{16+4+4+16}{10}=4$ and the standard deviation is $2$.
\end{itemize}
All four statistics equal $\boxed{2}$.

\textbf{How it was found.} Use a set symmetric about $m$, so the mean and median are $m$. With ten values $m\pm a,m\pm b,m\pm c,m\pm d,m\pm e$ ($0\le a\le b\le c\le d\le e$), the quartiles are $m\mp c$, so the IQR is $2c$, and the variance is $\dfrac{a^2+b^2+c^2+d^2+e^2}{5}$. Setting $m=2c$ and the variance equal to $m^2=4c^2$ gives $a^2+b^2+d^2+e^2=19c^2$; with $c=1$, $a=b=d=1$ and $e=4$ work. (The outlier $e$ is what makes the standard deviation as large as the IQR.)
